Orcus (2004 DW), un objecte del Sistema Solar descobert el febrer del 2004, és un dels cossos celestes més grans del cinturó de Kuiper i un dels candidats a ser considerat, en el futur, planeta nan per la Unió Astronòmica Internacional (UAI). Orcus té, aproximadament, una massa de $6,41\times 10^{20}\ \mathrm{kg}$, un radi de 459~km i un període orbital de 248 anys.

\figura[0.5]{fig1.png}

\begin{apartat}{1}
Calculeu la distància mitjana entre Orcus i el Sol en unitats astronòmiques (UA).
\end{apartat}

\begin{apartat}{1}
Determineu la velocitat d'escapament (deduïu la fórmula tenint en compte l'energia del cos que s'escapa) i la intensitat del camp gravitatori a la seva superfície.
\end{apartat}

\dades{%
  Radi orbital mitjà de la Terra $= 1,00$~UA.\\
  Període orbital de la Terra $= 1,00$~any.\\
  $G = 6,67\times 10^{-11}\ \mathrm{N\,m^2\,kg^{-2}}$.}
